\noindent Q : \textbf{1. Perusahaan ini bergerak dibidang apa ?}
\\ A : Perusahaan ini bergerak di bidang IT sebagai sistem integrator yang menjual barang-barang dari principal dan menyesuaikannya dengan nilai-nilai perusahaan.

\noindent Q : \textbf{2. Apakah perusahaan ini memiliki divisi yang melakukan kerja lembur ?}
\\ A : Ya, terdapat beberapa divisi seperti \textit{Infrastructure} dan \textit{Business Support} yang melakukan kerja lembur.

\noindent Q : \textbf{3. Apakah terdapat pencatatan kerja lembur staff ?}
\\ A : Ada, namun masih dilakukan secara manual menggunakan \textit{Microsoft Excel}. Meskipun ada sistem HRIS, sistem tersebut belum memadai untuk pencatatan lembur.

\noindent Q : \textbf{4. Bagaimana cara kerja lembur saat ini dicatat ?}
\\ A : Prosesnya dimulai dengan \textit{request} surat perintah lembur yang harus disetujui oleh atasan di akhir bulan. Proses ini memakan waktu lama, kurang efektif, dan memerlukan pengecekan \textit{manual} yang memakan banyak waktu.

\noindent Q : \textbf{5. Apa saja keluhan saat mencatat jam kerja lembur ?}
\\ A : Pencatatan \textit{manual} menyebabkan banyak kesalahan manusia (\textit{human error}). Kami ingin mempermudah langkah dan perhitungan lembur agar bisa langsung dari absensi, sehingga saat engineer \textit{clock in} dan \textit{clock out}, lemburnya langsung dihitung.

\noindent Q : \textbf{6. Dimana pelaksanaan lembur kerja biasanya dilakukan?}
\\ A : Kerja lembur bisa dilakukan secara \textit{remote}, sehingga penting memiliki sistem absensi yang memungkinkan \textit{clock in} dan \textit{clock out} dari mana saja.

\noindent Q : \textbf{7. Siapa saja yang bertanggung jawab untuk menyetujui dan memverifikasi data jam kerja lembur?}
\\ A : Proses persetujuan dilakukan secara bertahap mulai dari \textit{supervisor}, kemudian disetujui oleh \textit{manager}, dan terakhir masuk ke HR.

\noindent Q : \textbf{8. Jika terdapat aplikasi untuk mencatat jam lembur, fitur apa saja yang dibutuhkan ?}
\\ A : Aplikasi tersebut harus memiliki fitur absensi untuk \textit{clock in} dan \textit{clock out}, menghitung lembur berdasarkan jam, meminimalkan sistem administrasi, dan dapat menghitung lembur secara otomatis.

\noindent Q : \textbf{9. Apakah saat ini sudah ada sistem aplikasi presensi \textit{online}?}
\\ A : Ya, kami memiliki aplikasi bernama "Andal", tetapi aplikasi ini belum mampu mencakup semua kebutuhan.

\noindent Q : \textbf{10. Apakah aplikasi ini milik \textit{internal} perusahaan atau menggunakan pihak \textit{external} ? }
\\ A : Aplikasi ini dari pihak \textit{external} dengan sistem pembayaran sekali (\textit{one-time payment}).